% ═══════════════════════════════════════════════════════════════════════════════
%  CAPÍTULO — MEmento
% ═══════════════════════════════════════════════════════════════════════════════

\chapter{Memento}

% ─── Situación Problema ──────────────────────────────────────────────────────
\section{Situación Problema}

Cuando un usuario de Airbnb configura una búsqueda con múltiples filtros
—destino, fechas, tipo de alojamiento, amenidades— y luego navega hacia
la ficha de una propiedad, al presionar el botón «atrás» espera que su
búsqueda esté exactamente como la dejó: los mismos filtros activos, el
mismo punto de desplazamiento en la lista de resultados y la misma vista
seleccionada (mapa o lista). Esto mismo aplica cuando el usuario
interrumpe una reserva a mitad del proceso multistep y regresa más tarde.

Restablecer el estado de la búsqueda requiere conservar todos sus
atributos en un momento específico. Exponer el estado interno del objeto
\texttt{BusquedaActiva} para que capas externas lo almacenen viola el
encapsulamiento y acopla el mecanismo de almacenamiento a los detalles
internos del objeto.

% ─── Solución ────────────────────────────────────────────────────────────────
\section{Solución}

El patrón Memento permite que \texttt{BusquedaActiva} produzca un objeto
\texttt{InstantaneaBusqueda} que captura su estado interno completo sin
exponerlo al exterior. El \texttt{GestorNavegacion} almacena estas
instantáneas en una pila; cuando el usuario regresa, solicita la
instantánea correspondiente y \texttt{BusquedaActiva} se restaura al
estado guardado.

El encapsulamiento del objeto de búsqueda se mantiene intacto —el
gestor maneja instantáneas opacas—, y la experiencia del usuario es
fluida al retomar exactamente donde lo dejó, independientemente de cuántos
pasos de navegación hayan ocurrido en el intermedio.

% ─── Diagrama de clases ─────────────────────────────────────────────────────
\begin{figure}[H]
    \centering
    % \includegraphics[width=\textwidth]{imagenes/abstract_factory_diagrama.png}
    \caption{Diagrama de clases del patrón Abstract Factory}
\end{figure}


% ─── Roles de las Clases ────────────────────────────────────────────────────
\section{Roles de las Clases}

% Explicar la responsabilidad de cada clase/interfaz

\begin{itemize}
    \item \textbf{AbstractFactory:}
    \item \textbf{ConcreteFactory:}
    \item \textbf{AbstractProduct:}
    \item \textbf{ConcreteProduct:}
    \item \textbf{Client:}
\end{itemize}


% ─── Main ───────────────────────────────────────────────────────────────────
\section{Main}

% Explicar qué realiza el programa principal

\begin{lstlisting}[language=Java, caption=NombreClase]
 
\end{lstlisting}


% ─── Salida del Main ────────────────────────────────────────────────────────
\section{Salida del Main}

\begin{lstlisting}[language=Java, caption=NombreClase]
 
\end{lstlisting}


% ─── Clases ─────────────────────────────────────────────────────────────────
\section{Clases}

% ─── Clase 1 ────────────────────────────────────────────────────────────────
\subsection{NombreClase}

\begin{lstlisting}[language=Java, caption=NombreClase]
 
\end{lstlisting}


% ─── Clase 2 ────────────────────────────────────────────────────────────────
\subsection{NombreClase}

\begin{lstlisting}[language=Java, caption=NombreClase]
 
\end{lstlisting}


% ─── Clase 3 ────────────────────────────────────────────────────────────────
\subsection{NombreClase}

\begin{lstlisting}[language=Java, caption=NombreClase]
 
\end{lstlisting}


% ─── Conclusiones ───────────────────────────────────────────────────────────
\section{Conclusiones}

% Explicar:
% - Ventajas del patrón
% - Cuándo usarlo
% - Beneficios obtenidos
% - Posibles desventajas
% - Relación con principios SOLID